\documentclass[11pt]{article}
\usepackage{amsmath,amsthm,amssymb}
\usepackage[margin=1in]{geometry}
\usepackage{enumitem}

\newtheorem{theorem}{Theorem}
\newtheorem{lemma}[theorem]{Lemma}
\newtheorem{corollary}[theorem]{Corollary}
\newtheorem{proposition}[theorem]{Proposition}
\theoremstyle{definition}
\newtheorem{definition}[theorem]{Definition}
\newtheorem{remark}[theorem]{Remark}

\title{The $q$-integer $[3]_q$ as the engine of Theorem B,\\
and the polynomial-vs.-multiset-positivity gap}
\author{Clio}
\date{6 May 2026 (evening)}

\begin{document}
\maketitle

\begin{abstract}
The May~6 day-session produced five writeups establishing Theorems A, B, and C
at the multiset level, the path-level Theorem B with explicit $n=5$ injection,
the refutation of Open Question Q3, and the structural reduction of the
sign-positivity lemma to two single-component multiplicity bounds. This evening
note adds one cleanly proved observation that ties them together. The sole
``perturbation'' separating $\mathrm{atom}_{\mathrm{RTL}}$ from $A_{(3,2)}$ is the
$q$-integer multiple $2q[3]_q$, where $[3]_q := 1 + q + q^2$. Convolving this
into the $n=6$ identity gives
\[
A_{(3,2,1)}(q) \;=\; A_{(3,1,1)}(q)\, A_{(3,2)}(q) \;+\; 2q [3]_q\, A_{(3,1,1)}(q),
\]
a polynomial decomposition whose perturbation $2q[3]_q\, A_{(3,1,1)}$ is itself
palindromic, non-negative, and factors cleanly. The same perturbation, however,
has \emph{signed} canonical $\sigma_1$-G1 multiplicity $(0, +2, +2, -4)$ on $B_6$.
The polynomial is non-negative but not multiset-positive in the canonical basis.
This is exactly the obstruction to a direct multiset interpretation of the
perturbation, and is the precise structural origin of the ``hidden'' atom
$\mathrm{atom}_{\mathrm{RTL}}$. We take the result as a clean structural cap on
the day's program: the gap between polynomial- and multiset-positivity of the
perturbation $2q[3]_q\, A_{(3,1,1)}$ is what forces the rewrite into
$\mathrm{atom}_{\mathrm{RTL}}$, and is what makes Theorem B a genuine structural
statement rather than a polynomial identity.
\end{abstract}

\section{Setup}

We adopt the framework of \texttt{2026-05-06-twin-multiset.tex},
\texttt{2026-05-06-theorem-B-multiset.tex}, and
\texttt{2026-05-06-sign-positivity-structural.tex}. Recall the canonical
$\sigma_1$-G1 basis $B_{2k} = \{q^c (1+q)^{2k-2c} : c = 0,\dots,k\}$, which is
$\mathbb{Q}[q]$-linearly independent. Each cell-character $A_\lambda(q)$ admits
a unique non-negative integer expansion in $B_{2 m_\lambda}$, with multiplicity
vector $M_\lambda$, and multiplication of polynomials becomes convolution of
multiplicity vectors. The atoms relevant here:
\begin{equation}\label{eq:vectors}
M_{(3,1,1)} = (1,2),\quad M_{(3,2)} = (1,5,3),\quad M_{(4,1)} = (1,3,5),\quad
M_{(3,2,1)} = (1,9,15,2),\quad M_{(5,1)} = (1,3,9,14).
\end{equation}

The polynomial $[k]_q := 1 + q + \dots + q^{k-1}$ is the $q$-integer of order $k$.
We will use $[3]_q = 1 + q + q^2$ throughout.

\section{The $q$-integer identity}

\begin{lemma}[$2q$-difference, restated]\label{lem:2q}
\[
A_{(3,2)}(q) - A_{(4,1)}(q) \;=\; 2q \cdot [3]_q.
\]
\end{lemma}

\begin{proof}
By direct expansion of $M_{(3,2)} - M_{(4,1)} = (0, +2, -2)$ on $B_4$:
\[
2 q (1+q)^2 - 2 q^2 \;=\; 2q\bigl[(1+q)^2 - q\bigr] \;=\; 2q\bigl[1 + q + q^2\bigr] \;=\; 2q[3]_q.
\]
The middle equality is $(1+q)^2 = 1 + 2q + q^2 = [3]_q + q$.
\end{proof}

\begin{lemma}[Multiplicity vector of $[3]_q \cdot A_{(3,1,1)}$]\label{lem:3times311}
\[
[3]_q\, A_{(3,1,1)}(q) \;=\; 1 + 5q + 6q^2 + 5q^3 + q^4.
\]
The polynomial is palindromic of degree $4$, with all coefficients non-negative.
\end{lemma}

\begin{proof}
$A_{(3,1,1)}(q) = (1+q)^2 + 2q = 1 + 4q + q^2$. Then
\[
(1 + q + q^2)(1 + 4q + q^2) \;=\; 1 + (1+4)q + (1+4+1)q^2 + (4+1)q^3 + q^4 \;=\; 1 + 5q + 6q^2 + 5q^3 + q^4.
\]
Palindromicity is by inspection.
\end{proof}

\begin{theorem}[Polynomial decomposition of $A_{(3,2,1)}$]\label{thm:decomp}
\[
A_{(3,2,1)}(q) \;=\; A_{(3,1,1)}(q)\, A_{(3,2)}(q) \;+\; 2q\, [3]_q\, A_{(3,1,1)}(q),
\]
and equivalently
\[
A_{(5,1)}(q) \;=\; A_{(3,1,1)}(q)\, A_{(4,1)}(q) \;-\; 2q\, [3]_q\, A_{(3,1,1)}(q).
\]
The perturbation polynomial $2q\, [3]_q\, A_{(3,1,1)}(q) = 2(q + 5q^2 + 6q^3 + 5q^4 + q^5)$
is palindromic, non-negative, and supported on $\bar c = 1, \dots, 5$ (degree
$1$ through $5$ in $q$).
\end{theorem}

\begin{proof}
The multiset Theorem~B (Theorem~3.1 of \texttt{2026-05-06-theorem-B-multiset.tex}) is
$M_{(3,2,1)} = M_{(3,1,1)} \ast \mathrm{atom}_{\mathrm{RTL}}$ where
$\mathrm{atom}_{\mathrm{RTL}} = 2 M_{(3,2)} - M_{(4,1)}$. Translating to polynomials,
\[
A_{(3,2,1)} \;=\; A_{(3,1,1)} \cdot \bigl(2 A_{(3,2)} - A_{(4,1)}\bigr) \;=\; A_{(3,1,1)}\, A_{(3,2)} \;+\; A_{(3,1,1)}\bigl(A_{(3,2)} - A_{(4,1)}\bigr).
\]
By Lemma~\ref{lem:2q}, $A_{(3,2)} - A_{(4,1)} = 2q[3]_q$, so the second summand
is $A_{(3,1,1)} \cdot 2q[3]_q$. Multiplying out using
Lemma~\ref{lem:3times311}: $2q [3]_q\, A_{(3,1,1)}(q) = 2q (1 + 5q + 6q^2 + 5q^3 + q^4) = 2q + 10q^2 + 12q^3 + 10q^4 + 2q^5$, palindromic.

The $A_{(5,1)}$-form is analogous, using $A_{(5,1)} = A_{(3,1,1)}(2 A_{(4,1)} - A_{(3,2)})$ and the negative of the same identity.
\end{proof}

\begin{remark}[Numerical check]\label{rem:check}
Direct expansion:
\begin{align*}
A_{(3,1,1)}(q)\, A_{(3,2)}(q) &= (1 + 4q + q^2)(1 + 9q + 19q^2 + 9q^3 + q^4) \\
&= 1 + 13q + 56q^2 + 94q^3 + 56q^4 + 13q^5 + q^6,\\
2q [3]_q\, A_{(3,1,1)}(q) &= 2q + 10q^2 + 12q^3 + 10q^4 + 2q^5,\\
\text{sum} &= 1 + 15q + 66q^2 + 106q^3 + 66q^4 + 15q^5 + q^6 \;=\; A_{(3,2,1)}(q). \checkmark
\end{align*}
\end{remark}

\section{The polynomial-vs.-multiset positivity gap}

The polynomial $2q[3]_q\, A_{(3,1,1)}(q)$ has all non-negative coefficients (it is a
positive polynomial). But its expansion in the canonical $\sigma_1$-G1 basis
$B_6$ involves a negative entry. This is the structural origin of why
$\mathrm{atom}_{\mathrm{RTL}}$ exists as a hidden multiset rather than as a
direct ``$A_{(3,1,1)}\, p(q)$'' multiset for some positive multiset $p$.

\begin{lemma}[Signed $B_6$-decomposition of the perturbation]\label{lem:perturb-Bdec}
On $B_6 = \{(1+q)^6, q(1+q)^4, q^2(1+q)^2, q^3\}$, the perturbation has unique signed expansion
\[
2q [3]_q\, A_{(3,1,1)}(q) \;=\; 0 \cdot (1+q)^6 + 2 \cdot q(1+q)^4 + 2 \cdot q^2(1+q)^2 + (-4) \cdot q^3.
\]
That is, signed $B_6$-multiplicity $(0, +2, +2, -4)$.
\end{lemma}

\begin{proof}
By Theorem~\ref{thm:decomp}, $A_{(3,2,1)} - A_{(3,1,1)} A_{(3,2)} = 2q[3]_q\, A_{(3,1,1)}$,
which corresponds on $B_6$ to $M_{(3,2,1)} - M_{(3,1,1)} \ast M_{(3,2)} = (1,9,15,2) - (1,2)\ast(1,5,3) = (1,9,15,2) - (1, 7, 13, 6) = (0, +2, +2, -4)$.

Verifying the polynomial expansion: $0 + 2 q(1+q)^4 + 2 q^2 (1+q)^2 - 4 q^3$.
\begin{align*}
2 q (1+q)^4 &= 2(q + 4q^2 + 6q^3 + 4q^4 + q^5) = 2q + 8q^2 + 12 q^3 + 8q^4 + 2q^5,\\
2 q^2 (1+q)^2 &= 2(q^2 + 2q^3 + q^4) = 2q^2 + 4q^3 + 2q^4,\\
-4 q^3 &= -4 q^3.
\end{align*}
Sum: $2q + (8+2)q^2 + (12+4-4)q^3 + (8+2)q^4 + 2q^5 = 2q + 10q^2 + 12q^3 + 10q^4 + 2q^5$, matching Theorem~\ref{thm:decomp}. The expansion is unique by the $\mathbb{Q}[q]$-linear independence of $B_6$.
\end{proof}

\begin{corollary}[The polynomial is positive but not multiset-positive]\label{cor:gap}
The polynomial $2q[3]_q\, A_{(3,1,1)}(q)$ has all non-negative coefficients in $\mathbb{Z}[q]$, hence is positive in the polynomial sense. But its canonical $\sigma_1$-G1 multiplicity vector $(0, +2, +2, -4)$ on $B_6$ has a negative entry at $\bar c = 3$. So the perturbation is \emph{not} a non-negative element of the $\sigma_1$-G1 multiset cone over $B_6$.
\end{corollary}

\begin{proof}
Both statements are explicit: the polynomial coefficients $(0, 2, 10, 12, 10, 2, 0)$ are non-negative in $\mathbb{Z}[q]$; the $B_6$-multiplicity $(0, +2, +2, -4)$ has a negative third entry.
\end{proof}

\begin{remark}[Why this matters]\label{rem:matter}
This corollary is the structural reason why Theorem~B at the multiset level is non-trivial. If the perturbation $2q[3]_q\, A_{(3,1,1)}$ had non-negative $B_6$-multiplicity, then $M_{(3,2,1)}$ would simply be $M_{(3,1,1)} \ast M_{(3,2)}$ plus a non-negative multiset, with no need for the hidden atom rewrite $\mathrm{atom}_{\mathrm{RTL}} = 2 M_{(3,2)} - M_{(4,1)}$. But the perturbation has a $-4$ at $\bar c = 3$, so this naive ``add-positive'' reading fails. The rewrite $M_{(3,2,1)} = M_{(3,1,1)} \ast \mathrm{atom}_{\mathrm{RTL}}$ absorbs the perturbation by introducing $-M_{(4,1)}$ inside $\mathrm{atom}_{\mathrm{RTL}}$.

Concretely: at $\bar c = 3$, $M_{(3,1,1)} \ast M_{(3,2)} = 6$, but $M_{(3,2,1)} = 2$. We need to \emph{remove} $4$ paths-worth at this grade. In multiset language this is impossible (you cannot have a negative number of paths). The trick is to rewrite using the doubled atom: $M_{(3,1,1)} \ast 2 M_{(3,2)} = 12$, and $M_{(3,1,1)} \ast M_{(4,1)} = 10$, with difference $12 - 10 = 2$. This is the multiset Theorem B (or, the path-level Theorem B as a containment-and-complement).

So the gap between polynomial- and multiset-positivity, exhibited
quantitatively by Lemma~\ref{lem:perturb-Bdec}, is the precise structural
origin of $\mathrm{atom}_{\mathrm{RTL}}$ as a hidden multiset.
\end{remark}

\section{What the $[3]_q$ factor suggests}

The factorization in Theorem~\ref{thm:decomp} isolates the $q$-integer $[3]_q$
as the sole non-trivial factor in the perturbation. We make a few observations.

\subsection*{$[3]_q$ vs $A_{(2,1)}$ at $S_3$}

The polynomial $[3]_q = 1 + q + q^2$ is the $q$-character of the trivial
representation of $S_2$ induced to $S_3$, equivalently the $q$-Hilbert series of
projective space $\mathbb{P}^2$ (or of any flag variety with three Bruhat cells of consecutive lengths). The post-prefix
$\sigma_1$-G1 atom $A_{(2,1)}(q)$ at $S_3$ has the same degree (degree $2$) and is also
palindromic. We have not verified whether $A_{(2,1)}(q) = [3]_q$ at $S_3$ in this
note (the comparison requires recomputing $A_{(2,1)}$ in the $\sigma_1$-G1
post-prefix convention), but the coincidence is suggestive: if it holds, then
$2q [3]_q = 2 q\, A_{(2,1)}(q)$ and Theorem~\ref{thm:decomp} reads
\[
A_{(3,2,1)}(q) \;=\; A_{(3,1,1)}(q)\, A_{(3,2)}(q) \;+\; 2 q\, A_{(3,1,1)}(q)\, A_{(2,1)}(q).
\]
This would give the perturbation a clean partition-theoretic identity.

\subsection*{Open question: which other partition pairs admit a $q$-integer perturbation?}

The pair $((3,2), (4,1))$ at $S_5$ is special because $A_{(3,2)} - A_{(4,1)} = 2q[3]_q$ is a clean $q$-integer multiple, with $[3]_q$ palindromic of degree~$2$ ($n=3$ in the $q$-integer family). A natural question is: which other partition pairs $(\lambda, \mu)$ at $S_n$ have $A_\lambda - A_\mu = c \cdot q^a \cdot [k]_q$ for some positive integer $c$, non-negative integer $a$, and integer $k \geq 2$? Each such pair would give an analogous Theorem-B-like multiset rewrite. The next session should test this systematically at $n = 6, 7$ pairs.

\section{What this does and does not prove}

\textbf{What this proves.} The clean polynomial decomposition
$A_{(3,2,1)} = A_{(3,1,1)} A_{(3,2)} + 2q [3]_q A_{(3,1,1)}$, the corresponding
$A_{(5,1)} = A_{(3,1,1)} A_{(4,1)} - 2q [3]_q A_{(3,1,1)}$, and the explicit
signed $B_6$-multiplicity vector $(0, +2, +2, -4)$ of the perturbation. The
polynomial is non-negative but not multiset-non-negative.

\textbf{What this does not prove.} It does \emph{not} construct a structural
W-graph-natural injection of paths, does not prove canonicity, and does not
prove either of the two open empirical identities $D(q) = \mathrm{cross}_{(3,1,1)\to(3,2)}(q)$
(Theorem~4 of \texttt{2026-05-06-paths-explicit-and-Q3-refuted.tex}) or any
generalization beyond the specific $n=5,6$ pair.

\textbf{Status of the May 6 program at end of evening.} The polynomial identities
of Theorems A, B, C are all verified. Their multiset upgrades are all proved
(Theorem~B with hidden $\mathrm{atom}_{\mathrm{RTL}} = (1, 7, 1)$ and
$\mathrm{atom}_{\mathrm{LTR}} = (1, 1, 7)$ as the structural inner factors). The
path-level Theorem~B is proved as an existence-of-injection statement (with an
explicitly tabulated lex injection). Open Question Q3 of
\texttt{2026-05-06-theorem-B-multiset.tex} is refuted; in its place we have the
empirical identity $D(q) = \mathrm{cross}_{(3,1,1)\to(3,2)}(q)$ from
\texttt{2026-05-06-paths-explicit-and-Q3-refuted.tex}, whose structural meaning
is open. The sign-positivity lemma is proved and its $W$-graph-natural status
is reduced to two single-component multiplicity bounds (Theorem~1 of
\texttt{2026-05-06-sign-positivity-structural.tex}). The current note adds the
$[3]_q$-isolation of the perturbation and the polynomial-vs.-multiset
positivity gap as the precise structural fact behind the day's program.

\section*{Computational verification}

\begin{itemize}[leftmargin=2em,topsep=0.2em,itemsep=0.1em]
\item $A_{(3,2)}(q) - A_{(4,1)}(q) = 2q(q^2+q+1) = 2q[3]_q$
  (April~29 polynomial residual zero;
  Lemma~3 of \texttt{2026-05-06-theorem-B-multiset.tex}; Lemma~\ref{lem:2q} above).
\item $[3]_q\, A_{(3,1,1)}(q) = 1 + 5q + 6q^2 + 5q^3 + q^4$
  (direct expansion; Lemma~\ref{lem:3times311}).
\item $A_{(3,1,1)}(q)\, A_{(3,2)}(q) = 1 + 13q + 56q^2 + 94q^3 + 56q^4 + 13q^5 + q^6$
  (Remark~\ref{rem:check}).
\item $A_{(3,2,1)}(q) = 1 + 15q + 66q^2 + 106q^3 + 66q^4 + 15q^5 + q^6$
  (\texttt{2026-05-06-twin-multiset.tex}).
\item Difference $A_{(3,2,1)} - A_{(3,1,1)} A_{(3,2)} = 2q + 10q^2 + 12q^3 + 10q^4 + 2q^5 = 2q[3]_q\, A_{(3,1,1)}$ (Theorem~\ref{thm:decomp}).
\item $B_6$-decomposition of the perturbation: $(0, +2, +2, -4)$ (Lemma~\ref{lem:perturb-Bdec}).
\end{itemize}

\end{document}
